\documentclass[10pt,letterpaper]{article}
% Spanish CV. Build with: ./cv/build.sh  (layout lives in style.tex)
% Shared layout for CV.tex (English) and CV_ES.tex (Spanish).
% Moved verbatim from the Overleaf preamble; change styling here, once, for both languages.

% ---------- Packages ----------
\usepackage[margin=0.6in]{geometry}
\usepackage[hidelinks]{hyperref}
\usepackage{titlesec}
\usepackage{enumitem}
\usepackage[T1]{fontenc}
\usepackage{tgtermes}
\usepackage{microtype}

% ---------- Formatting ----------
\pagestyle{empty}
\setlength{\parindent}{0pt}
\setlength{\parskip}{0pt}

\titleformat{\section}
  {\bfseries\uppercase}
  {}
  {0pt}
  {}
  [\vspace{-2pt}\titlerule]

\titlespacing{\section}{0pt}{7pt}{4pt}

\setlist[itemize]{
  leftmargin=12pt,
  topsep=1pt,
  itemsep=2pt,
  parsep=0pt
}

% \entry{Organization}{Dates}{Role or degree}{Location}
% The trailing % signs stop line ends from adding stray spaces (which pushed the
% right-aligned location of a row ending a paragraph 2.5pt short of the margin).
\newcommand{\entry}[4]{%
  \vspace{2pt}%
  \textbf{#1} \hfill #2 \\
  \textit{#3} \hfill \textit{#4}%
}


\begin{document}

% ---------- Encabezado ----------
\begin{center}
    {\LARGE \textbf{Carlos Alberto Bustamante Gaytán}} \\[-1pt]
    Ingeniero Cloud \& MLOps | AWS | Kubernetes | Terraform | Infraestructura de IA \\[2pt]
    Estado de México, México |
    \href{mailto:me@carlosbustamante.dev}{me@carlosbustamante.dev} |
    +52 55 3131 1250 \\
    \href{https://carlosbustamante.dev}{carlosbustamante.dev} |
    \href{https://github.com/charlybgai}{github.com/charlybgai} |
    \href{https://linkedin.com/in/charlybgai}{linkedin.com/in/charlybgai}
\end{center}

% ---------- Resumen ----------
\section*{Resumen Profesional}
Ingeniero Cloud y MLOps con trayectoria en el diseño y despliegue de sistemas de IA en producción sobre AWS y Kubernetes. Sólida experiencia en Terraform, Python, despliegues en contenedores, model serving, infraestructura para LLMs y sistemas backend escalables. Enfocado en plataformas cloud confiables para automatización con agentes de IA y cargas de ML en producción.

% ---------- Experiencia ----------
\section*{Experiencia Profesional}

\entry
{GlobalLogic (cliente: Ring, empresa de Amazon)}{Junio 2026 -- Presente}
{Senior Software Engineer}{Remoto}
\begin{itemize}
  \item Desarrollo de agentes de IA que agilizan flujos de trabajo internos de producto e ingeniería en Ring para PMs, TPMs e ingenieros.
\end{itemize}

\entry
{Corteza.ai}{Enero 2026 -- Mayo 2026}
{Senior Machine Learning Engineer}{Remoto}
\begin{itemize}
  \item Diseño de la infraestructura de IA en producción sobre AWS con Terraform y Kubernetes, para cargas de trabajo de ML y LLMs en EKS, ECR, EC2, S3, EFS y Route 53.
  \item Creación de AMIs administradas con Terraform que redujeron el tiempo de arranque de pods en Kubernetes de $\sim$5 a 2 minutos, al minimizar la descarga de imágenes de contenedor.
  \item Migración de un asistente en producción, de procesamiento con RabbitMQ a inferencia en streaming con vLLM, lo que aumentó las respuestas concurrentes de 2 a 10 en una sola instancia con GPU.
\end{itemize}

\entry
{Corteza.ai}{Septiembre 2024 -- Enero 2026}
{Machine Learning Engineer}{Remoto}
\begin{itemize}
  \item Liderazgo de flujos completos de fine-tuning para LLMs internos: extracción de código fuente, creación de datasets, preparación del entrenamiento, fine-tuning supervisado y evaluación.
  \item Desarrollo en Python de inferencia para LLMs, modelos de detección de objetos y de difusión, y APIs de IA generativa de terceros.
  \item Creación y despliegue de servicios de IA en contenedores para aplicaciones en producción y herramientas internas de ML.
\end{itemize}

\entry
{Tecnológico de Monterrey}{Enero 2024 -- Junio 2024}
{Profesor de Cátedra}{Estado de México, México}
\begin{itemize}
  \item Docencia en programación, desarrollo de aplicaciones y electrónica; mentoría en proyectos prácticos con micro:bit.
\end{itemize}

\entry
{TELUS International}{Febrero 2020 -- Junio 2022}
{Data Analyst}{Remoto}
\begin{itemize}
  \item Análisis y anotación de contenido digital para mejorar la calidad de los datos de entrenamiento de IA, búsqueda y recomendación.
\end{itemize}

% ---------- Habilidades ----------
\section*{Habilidades}
\textbf{Cloud e infraestructura:} AWS, Azure, Linux, EKS, ECR, EC2, S3, EFS, Route 53, CloudWatch \\
\textbf{DevOps e IaC:} Terraform, AWS CDK, Kubernetes, Docker, Git, CI/CD, despliegues en contenedores \\
\textbf{IA/ML y MLOps:} PyTorch, TensorFlow, Scikit-learn, LLMs, vLLM, entrenamiento, fine-tuning, serving y despliegue de modelos \\
\textbf{Programación y datos:} Python, TypeScript, FastAPI, REST APIs, SQL, Pandas, NumPy, PySpark \\
\textbf{Idiomas:} español (nativo), inglés (avanzado), alemán (intermedio)

% ---------- Certificaciones ----------
\section*{Certificaciones}
\begin{itemize}
  \item AWS Certified Generative AI Developer -- Professional \hfill Vigente hasta 2029
  \item AWS Certified Machine Learning Engineer -- Associate \hfill Vigente hasta 2029
  \item AWS Certified Solutions Architect -- Associate \hfill Vigente hasta 2027
  \item HashiCorp Certified: Terraform Associate (004) \hfill Vigente hasta 2028
  \item Microsoft Certified: Azure Data Scientist Associate \hfill Vigente hasta 2027
\end{itemize}

% ---------- Educación ----------
\section*{Educación}

\entry
{Tecnológico de Monterrey}{Agosto 2022 -- Junio 2024}
{Maestría en Ciencias de la Ingeniería}{Estado de México, México}
\begin{itemize}
  \item Coautoría de un artículo en IEEE (2023): controlador neuronal NARMA-L2 para la enseñanza de control basado en datos.
\end{itemize}

\entry
{Tecnológico de Monterrey}{Agosto 2017 -- Diciembre 2021}
{Ingeniería en Mecatrónica}{Estado de México, México}

\end{document}
